\documentclass[11pt,fleqn]{article}
\usepackage[margin=0.8in]{geometry}
\usepackage{amsmath,amssymb}
\usepackage[hidelinks]{hyperref}
\setlength{\mathindent}{0pt}

\newcommand{\X}{\mathcal X}
\newcommand{\C}{\mathcal C}
\newcommand{\Hy}{\mathcal H}
\newcommand{\err}{\operatorname{error}}
\newcommand{\eps}{\epsilon}
\newcommand{\Adv}{\operatorname{Adv}}

\begin{document}
\raggedright
\setlength{\parskip}{0.6em}

{\LARGE\bfseries Page 7: Implications and Further Reading\par}
\medskip

\section{What the Project Shows}

A major question we asked was whether having enough examples means that we are able to learn efficiently. For an $n$-bit PRF family, the finite-class bound gives us a sufficient sample size linear in $n$ for fixed accuracy and confidence. If we know the key, we can also evaluate the target efficiently. However, under PRF security, no polynomial-time learner can reliably turn that number of examples into an accurate predictor.

Even though this may at first seem off, these facts are compatible because they answer different questions. The sample bounds gives a number that tells us that any candidate predictor that is consistent with our examples is likely to be a good predictor. However, actually turning the examples into a predictor is difficult. We could try all $2^n$ keys and figure it out, but this is not efficient. We showed that there is no efficient algorithm to find this predictor, assuming PRF security holds.

\section{What the Hardness Claim Depends On}

The reduction we showed proves that the existence of an efficient learner would lead to the existence of an efficient distinguisher with non-negligible advantage. The conclusion is conditional on PRF security~\cite{mit}.

The learning guarantee from the main theorem is also over a uniform input distribution. Note that a distribution-free model would have to work on the uniform input distribution. An algorithm that could succeed on a specific (i.e. easier) distribution is making a different guarantee.

\section{Connection to Boolean Circuits}

Note that the PRF has a polynomial-time evaluable algorithm that takes in $k$ and $x$ as inputs. Thus, using our efficient TM-to-Circuit theorem from class, we could construct a circuit for each evaluator $F_k$~\cite{circuits}. We could hardwire the key bits and then leave the $n$ input bits free.

Thus, there is a polynomial size bound $p(n)$ large enough to contain a circuit for each member of the PRF family. A learner that could work for all functions represented by circuits of that size would then have to be able to learn PRFs. Thus, the same security assumption would give a learning-hardness consequence for a suitably size-bounded circuit class.

\section{Restricted Classes and Related Work}

For more restrictive target classes, we need new arguments that show that the hard functions are representable inside them. Kearns and Valiant give arguments for classes that include polynomial-size Boolean formulas, acyclic DFAs, and constant-depth threshold circuits. These reductions connect distribution-free learning for these classes to cryptographic and number-theoretic problems like RSA inversion~\cite{valiant}.

\emph{Further reading: Inherent Unpredictability}~\cite{vazirani}.

\section{Connection to Machine Learning}

When a problem seems hard, it may be useful to step back and question whether the limitation is coming from a lack of data, from a lack of training time, or from a breakdown in the assumptions the algorithm can make about the target and distribution.

\emph{Further reading: PAC Verification of Statistical Algorithms}~\cite{verification}.

\begin{thebibliography}{9}
\bibitem{mit}
MIT 6.S976 / 18.S996.
\emph{Cryptography and Machine Learning: Foundations and Frontiers}.
\href{https://mlcrypto.mit.edu/course/slides/lec5.pdf#page=47}
{Lecture 5, PRFs and learning hardness}, Spring 2026.

\bibitem{circuits}
An{\i}l Ada, with notes by Keming Meng and Evan Ning.
\emph{Boolean Circuits 2}. CS455, Fall 2026.

\bibitem{valiant}
Michael Kearns and Leslie Valiant.
\emph{Cryptographic Limitations on Learning Boolean Formulae and
Finite Automata}. Journal of the ACM, 41(1), 67--95, 1994.
\href{https://www.cis.upenn.edu/~mkearns/papers/cryptojacm.pdf}
{Paper}.

\bibitem{vazirani}
Michael J. Kearns and Umesh V. Vazirani.
\emph{An Introduction to Computational Learning Theory}.
MIT Press, 1994. Chapter 6, ``Inherent Unpredictability.''

\bibitem{verification}
Saachi Mutreja and Jonathan Shafer.
\emph{PAC Verification of Statistical Algorithms}.
Conference on Learning Theory, 2023.
\href{https://proceedings.mlr.press/v195/mutreja23a.html}{Paper}.
\end{thebibliography}

\end{document}